\documentclass[10pt]{article}
\usepackage[utf8]{inputenc}
\usepackage[T1]{fontenc}
\usepackage{amsmath}
\usepackage{amsfonts}
\usepackage{amssymb}
\usepackage[version=4]{mhchem}
\usepackage{stmaryrd}
\usepackage{graphicx}
\usepackage[export]{adjustbox}
\graphicspath{ {./images/} }
\usepackage{mathrsfs}
\usepackage{bbold}

\begin{document}
венства в (1) и (2) имеют место только в том случае, когда $w=f(z)$ есть дробно-линейная функция, отображающая круг $\Omega$ на себя.

Н е е в к л п д о в о, илиг и пе р болическ о е, ра с с т о я в и $d\left(z_{1}, z_{2}\right)$ есть расстояние в геометрии Лобачевского между точками $z_{1}$ и $z_{2}$, когда круг $\Omega$ принимается в качестве плоскости Лобачевского, а прямыми Лобачевского служат дуги окружностей, ортогональных единичной окружности (м о д е л ь ПІ уа нк а р е); при әтом

\[
\begin{gathered}
d\left(z_{1}, z_{2}\right)=\frac{1}{2} \ln \left(\alpha, \beta, z_{1}, z_{2}\right)= \\
=\frac{1}{2} \ln \frac{\left|1-\bar{z}_{1} z_{2}\right|+\left|z_{1}-z_{2}\right|}{\left|1-\bar{z}_{1} z_{2}\right|-\left|z_{1}-z_{2}\right|},
\end{gathered}
\]

где $\left(\alpha, \beta, z_{1}, z_{2}\right)$ - двойное отношение точек $z_{1}, z_{2}$ и определяемых әтими точками точек пересечения $\alpha, \beta$

\includegraphics[max width=\textwidth, center]{2023_07_08_4909cfdb4fc37cd3bb88g-1}
прямой Лобачевского, проходящей через $z_{1}$ и $z_{2}$, с единичной окружностью (см. рис.).

Неевклидова длина образа $f(L)$ любой спрямляемой кривой $L \subset \Omega$ при отображении $w=f(z)$ не превосходит неевклидовой длины $L$.

П. т. была установлена $\Gamma$. Пиком (см. [1]); дальнейпим ее обобшением является гиперболической метрики nринцип. В геометрич. теории функций из этих теорем выводятся оценки различных функционалов, связанных с отображаюпей функцией (см. [2], [3]).

Jum.: [1] P i c k G., "Math. Ann.», 1916, Bd 77, S. 1-6; [2] $\boldsymbol{\Gamma}$ о л у з и н $\Gamma$. М., Геометрическая теория функций комплексного переменного, 2 изд., М., 1966; [3] К а р а т е о д о р и К., Конформное отображение, пер. с англ., М.- Ј., 1934.

ПИКАРА ГРУППА - группа классов обратимых пучков (или линейных расслоений). Более точно, пусть $(X, \widehat{\widehat{O}} \mathrm{x})$ - окольцованное пространство. Пучок $\widehat{O}^{-}$-модулей $\mathscr{L}$ наз. о б р а т и м м м, если он локально изоморфен структурному пучку $\sigma_{X}$. Множество классов изоморфных обратимых пучков на $X$ обозначается Рic $(X)$. Тензорное произведение $\mathscr{L} \otimes$ $\otimes \sigma_{X} \mathscr{L}^{\prime}$ определяет на множестве Pic $(X)$ операцию, превращающую его в абелеву группу, наз. г р у п по й П и к а р а пространства $X$. Группа $\mathrm{Pic}(X)$ естественно изоморфна группе когомологий $H^{1}\left(X, \widehat{G}_{X}^{*}\right)$, где $\widehat{\sigma}_{X}^{*}$ - пучок обратимых элементов в $G x$.

Для коммутативного кольца $A$ группой Пикара $\mathrm{Pic} A$ наз. группа классов обратимых $A$-модулей; $\operatorname{Pic} A \cong \operatorname{Pic}(\operatorname{Spec} A)$. Для кольца Крулля $A$ группа $\mathrm{Pic} A$ тесно связана с классов дивизоров әруппой әтого кольца.

П. г. полного нормального алгебраич. многообразия обладает естественной алгебраич. структурой (см. Пикара схема). Связная компонента нуля группы $\mathrm{Pic}(X)$ обозначается $\mathrm{Pic}^{0}(X)$ и наз. многообразием Пикара для $X$; әто - алгебраич. группа (абелево многообразие, если многообразие $X$ гладко). Факторгруппа $\operatorname{Pic}(X) / \operatorname{Pic}^{0}(X)$ наз. г р у п п о й $\mathrm{H}$ е р о н а $-\mathrm{C}$ ев е р и и имеет конечное число образующих; ранг ее наз. ч и с л о м П и к а р а. В комплексном случае, когда $X$ - гладкое проективное многообразие над $\mathbb{C}$, группа $\mathrm{Pic}^{0}(X)$ изоморфиа факторгруппе пространства $H^{0}\left(X, \Omega_{X}^{1}\right)$ голоморфных 1-форм на $X$ по решетке $H^{1}(X, \mathbb{Z})$

Лит.: [1] М а м ф о р д Д., Лекции о кривых на алгебраической поверхности, пер. с англ., М., 1968. В. И. Данилов.

ПИКАРА МНОГООБРАЗИЕ полного гладкого алгебраического многообразия $X$ над алгебраически замкнутым полем - абелево многообразие $\mathfrak{B}(X)$, параметривующее факторгруппу $\operatorname{Div}^{a}(X) / P(X)$ группы $\operatorname{Div}^{a}(X)$ дивизоров, алгебраически әквивалентных нулю, по группе главных дивизоров $P(\mathrm{X})$, т. е. дивизоров рациональных функций. $\mathrm{C}$ точки зрения теории пучков II. м. параметризует множество классов изоморфных обратимых пучков с нулевым классом Чжәня, т. е. $\mathfrak{P}(X)$ совпадает со связной компонентой единицы $\mathrm{Pic}^{0}(X)$ Пикара epynпы $\operatorname{Pic}(X)$ многообразия $X$. Структура абелева многообразия на группе $\mathfrak{P}(X)=$ - $\operatorname{Div}^{a}(X) / P(X)$ однозначно характеризуется следующим свойством: для любого алгебраич. семейства дивизоров $D$ на $X$ с базой $S$ существует такое регулярное отображение $\varphi: S \rightarrow \mathfrak{B}(X)$, для к-рого $D(s)-$ $-D\left(s_{0}\right) \in \varphi(s)$, где $s_{0}$ - нек-рая фиксированная точка из $S_{0}$ (см. [2]). Размерность $q=\operatorname{dim} \Re(X)$ наз. и р р е г ул я $\mathrm{p}$ н о с т ь ю многообразия $X$.

Классической пример П. м.- Якоби многообразие гладкой проективной кривой. Другим примером может служить двойственное абелево многообразие (cM. [3]).

В случае, когда $X$ - гладкое проективное комплексное многообразие, $\mathfrak{B}(X)$ можно отождествить с группой обратимых аналитич. пучков на $X$, имеющих нулевой класс Чжәня [4]. Более того, многообразие Пикара $\mathfrak{B}(X)$ в этом случае изоморфно факторгруппе пространства $H^{1}\left(X, \widehat{G}_{X}\right)$ по рештетке $H^{1}\left(X, \mathbb{Z}_{3}\right) \subset H^{1}(X$, $\left.\sigma_{X}\right)$. В частности, иррегулярность $q$ многообразия $X$ совпадает с $\operatorname{dim} H^{1}\left(X, \hat{G}_{X}\right)=\operatorname{dim} H^{0}\left(X, \Omega_{X}^{1}\right)$, где $\Omega_{X}^{1}-$ пучок регулярных 1-форм. Последний результат верен также в случае неособых проективных кривых над любым алгебраически замкнутым полем и в случае полных гладких многообразий над алгебраически замкнутым полем характеристики 0 ; для произвольной характеристики имеет место лишь н е р а в е н с т в о И г у з ы $: \operatorname{dim} H^{1}\left(X, \widehat{O}_{X}\right) \geqslant q$ (известен пример гладкой алгебраич. поверхности $F$ иррегулярности 1 с $\operatorname{dim} H^{1}\left(F, G_{F}\right)=2$; см. [6]). Как видно отсюда, П.м. тесно связано с теорией одномерных дифференциальных форм. Начало исследования таких форм на римановых поверхностях положил Э. Пикар [1]; он доказал конечномерность пространства $H^{0}\left(X, \Omega_{X}^{1}\right)$ всюду регулярных форм.

Понятие П. м. может быть обобщено на случай полного нормального многообразия $X$. Изучалось также многообразие Пикара $\mathfrak{B}_{c}(X)$, соответствуюгее дивизорам Картье и обладающее, в отличие от $\mathfrak{P}(X)$, хорошими функториальными свойствами [9]. Многообразие $\mathfrak{H}_{c}(X)$ было построено для полных нормальных многообразий $X$ (см. [5]), а также для произвольных проективных многообразий [8].

Jum.: [1] P i c a r d E., "C. r. Acad. Sci.», 1884, t. 99, p. 961963; [2] III а Ф а р е в и ч И. Р., Основы алгебраической геометрии; М., 1972; [3] М а М Ф о р д Д., Абелевы многообразия, пер. с англ., М., 1971 ; [4] $\Gamma$ p и фф и т с Ф., Х а p р и с Д. Ж., М., 1982; [5] C h e v a 11 е у С., "Amer. J. Мath.", 1960, v. 82, p. 435-90; [6] I g u s a J.-I. "Proc. Nat. Acad. Sci. USA", 1955 , v. 41, № 11, p. 964-67; [7] M a ts u 8 a k a T., "Jap. J. Math.",
1951, v. $21, \mathcal{N}_{2}$, p. $217-36 ;[8]$ S e s h a d r C., "Ann. mat. pura ed appl.», 1962, v. 57, p. 117-42; [9] e r o " A e "Math. Ann.», 1965, Bd 158, № 3, S. 293-96.

ПИКАРА СХЕМА - естественное обобщение в рамках теории схем понятия Пикара многообразия $\mathfrak{B}(X)$ гладкого алгебраич. многообразия $X$. Для определения П. с. произвольной $S$-схемы $X$ рассматривается о т н осите льны й фу ункто категории $\operatorname{Sch} / S$ схем над схемой $S$. Значение этого функтора на $S$-схеме $S^{\prime}$ есть группа

\[
H^{0}\left(S^{\prime}, R_{f p q c}^{1} f_{*}^{\prime}\left(G_{m, X^{\prime}}\right)\right),
\]

где $f^{\prime}: X \times{ }_{s} S^{\prime} \rightarrow S^{\prime}$ есть морфизм замены базы, $R_{f p q c f^{\prime}}^{1}\left(G_{m, X^{\prime}}\right)$ - пучок ва топологии Гротендика


\end{document}